\chapter{Fracture Reduction}

\section*{What Is This Test or Procedure?}
Fracture reduction aligns displaced bone fragments to their normal anatomical position, promoting proper healing and restoring limb function.

\section*{Alternative Names}
Bone setting, Closed reduction, Open Reduction Internal Fixation (ORIF)

\section*{Principle}
Displaced bone fragments are realigned manually, by traction, or surgically so the fracture surfaces return to anatomic or near-anatomic apposition. Immobilization then holds the fragments in alignment while healing by callus formation proceeds.
\section*{History}
Hippocrates described traction and splinting techniques. Wilhelm Roentgen's discovery of X-rays in 1895 allowed visualization of alignment. Lister introduced aseptic open reduction.

\section*{Why This Test or Procedure Is Ordered}
Ordered for displaced bone fractures to prevent deformity, relieve pain, and avoid nerve or blood vessel compression.

\section*{Is This a Primary or Secondary Test or Procedure}
Primary definitive treatment for displaced limb fractures.

\section*{How This Test or Procedure Is Performed}
Closed reduction involves manipulating the bones externally under sedation or local block. Open reduction (ORIF) is a surgical procedure where bones are aligned directly and secured with plates, screws, or rods.

\section*{What Preparation Is Needed}
Perform neurovascular assessment of the distal limb. Obtain pre-reduction X-rays. Fast if sedation or general anesthesia is required.

\section*{Contraindications and Precautions}
None when a fracture is present. Active soft-tissue infection is a relative contraindication for immediate internal fixation.

\section*{Risks}
Nerve or vascular injury during manipulation, compartment syndrome, infection (in open reduction), or nonunion/malunion.

\section*{What Happens After the Test or Procedure}
Immobilize the limb immediately with a splint or cast. Obtain post-reduction X-rays to verify alignment. Re-assess distal pulses and sensation.

\section*{Understanding Your Results}
Post-reduction X-rays must show acceptable bone alignment and length restoration.

\section*{Normal Results}
Anatomical or acceptable functional alignment of bone fragments; intact distal neurovascular status.

\section*{Follow-up Tests and Procedures}
Follow-up X-rays in 1, 2, and 6 weeks to monitor healing; physical therapy.

\section*{References}
\begin{enumerate}
  \item MedlinePlus. Fracture Reduction. U.S. National Library of Medicine; 2025.
  \item Current Medical Diagnosis and Treatment. McGraw-Hill; 2025.
\end{enumerate}

\clearpage
